\section{证据规则}
\begin{enumerate}
  \item 调整证据能力 \yjt 能否作为证据。
        \begin{enumerate}
          \item 非法证据排除。
          \item 自白任意。不能强迫任何人\zd{自证其罪}。52 条。
          \item 最佳证据。能原件原件。
          \item 意见证据。约束\zd{证人猜测性、评论性、推断性的证言}。这么记，意谐\zd{疑}。
          \item 传闻证据。\bs{证人在法庭外}的陈述、\bs{非目击者}的当庭陈述（道听途说）。\zd{原则上不能}（立法没有落地）。
        \end{enumerate}
  \item 调整证明力 \yjt 证明作用大小。
        \begin{enumerate}
          \item \zd{补强证据}规则。证明力显然薄弱的证据 \yjt 主证据。题目出现\zd{印证}。
                \begin{enumerate}
                  \item 证人认知和表达\zd{存在一定困难}，但\zd{尚未丧失}正确认知、表达能力。
                  \item \zd{立场有问题、虚假动机}。有利的人作有利的证，不利的人作不利的证。
                  \item 要求。\bs{补强证据必须具有独立来源、必须具有证据能力、能够担保补强对象真实}。\al{测谎 \cuo}。
                \end{enumerate}
          \item 关联性规则。\al{类似事件、人的品格、表情、被害人过去行为 \cuo}。
        \end{enumerate}
  \item 非法证据排除。排非规定。
        \begin{enumerate}
          \item 哪些需要排除。
                \begin{enumerate}
                  \item \bs{言词排非暴、限、胁}。暴力、非法限制自由、胁迫。
                        \begin{enumerate}
                          \item 对被告人胁迫。侵犯本人、近亲属\zd{合法权益}。\al{让你儿子上不了考场 \dui，找你儿子取证 \cuo}。
                          \item 对证人、被害人一律排。
                          \item \bs{重复供述一并排，除非换人换阶段}。
                        \end{enumerate}
                  \item \bs{非法实物可补正}（相对客观）。应当予以补正 \hz 做出合理解释。毒树之果，洗洗能吃。
                \end{enumerate}
          \item 举证责任要求。谁来证明证据非法/合法。
                \begin{enumerate}
                  \item \zd{申请方。初步}证明 \yjt 合法性有疑问的程度。
                  \item \zd{检}察院。证明责任 \yjt 达到\zd{确实、充分}的标准。
                  \item \zd{法院要当庭作出决定}。
                \end{enumerate}
          \item 依申请（当法诉辩）\hz 依职权。
          \item 阶段和要求。\yi 侦查 \yjt 排非 14。\er 审查 \yjt 排非 17 \san 审判 \yjt 排非 23。\bs{都是应当}。
                \begin{enumerate}
                  \item 开庭前。
                        \begin{enumerate}
                          \item \zd{送达起诉书副本时}，应当告知当事人有权申请排非。
                          \item 原则\zd{开庭前}提出申请，\zd{庭审期间}发现相关线索、材料的除外。
                          \item \zd{被告及其辩护申请排非需要提供相关线索或材料}。
                          \item 应当\zd{召开庭前会议} \yjt 合法性\zd{有疑问的} \yjt 应当在\zd{庭审中调查}。
                        \end{enumerate}
                  \item 庭审中。
                        \begin{enumerate}
                          \item 原则应当\zd{先行当庭调查，最迟法庭调查结束前}完成证据的合法性调查。
                          \item 申请应当说明理由。确有疑问的 \yjt 应当进行调查。没有 \yjt 应当驳回申请。
                        \end{enumerate}
                  \item 二审。一审申请\zd{未审查}、一审审查\zd{我不服}、一审之后\zd{才发现}。
                \end{enumerate}
          \item 流程图。
        \end{enumerate}
\end{enumerate}
